% =====================================================================
%  Supplementary Material S1 -- The Reaction Admissibility Graph:
%  lineage, graph software, three-layer architecture, solver coupling
%
%  Drop-in replacement for the current S1.  Requires in the preamble:
%    \usepackage{amsmath,amssymb,booktabs,tabularx,xcolor,tikz}
%    \usetikzlibrary{arrows.meta,positioning,fit,backgrounds,calc,
%                    decorations.pathreplacing,patterns}
%  Bibliography keys are in S1_references.bib (merge into the main .bib).
%  Code paths refer to RadCluster_2_1 (https://github.com/Ghoniem/RadCluster).
% =====================================================================

\section{The Reaction Admissibility Graph: Lineage, Implementation, and Solver Coupling}
\label{sup:S1}
\begingroup
\setlength{\emergencystretch}{3em}% long code identifiers: avoid overfull lines

This section relates the reaction admissibility graph (RAG) of Sec.~2 to
chemical reaction network theory, describes its software representation
and the three-layer architecture of RadCluster, and explains how the graph
constrains the choice of integrator. A cluster vertex is
$\sigma=(\chi,n,p,\mathbf{c})$ (polarity, size, population, composition),
the graph is $\mathcal{G}=(\mathcal{V},\mathcal{E},\nu,k)$, and the master
equation is
\begin{equation}
  \frac{dc_a}{dt} \;=\; G_a \;+\; \sum_{r\in\mathcal{E}} S_{ar}\,J_r(\mathbf{c})
  \;-\; \mathcal{D}_a\,c_a ,
  \label{eq:S1_walker}
\end{equation}
with $S_{ar}$ the stoichiometric coefficient of vertex $a$ in edge $r$,
$J_r$ the flux built from kernel $k_r$, $G_a$ the cascade source and
$\mathcal{D}_a$ the loss to unresolved fixed sinks.

% ---------------------------------------------------------------------
\paragraph{Relation to chemical reaction network theory.}
Under mass-action kinetics a chemical reaction network (CRN) evolves as
$\dot{\mathbf{x}}=\mathsf{S}\,\mathbf{J}(\mathbf{x})$ with a constant
stoichiometric matrix $\mathsf{S}$
\cite{HornJackson1972,Feinberg1987,Feinberg2019,ErdiToth1989}.
Equation~\eqref{eq:S1_walker} has this form: vertices are species, edges
are reactions, binary edges (coalescence, annihilation) are directed
hyperedges with two tails \cite{Klamt2009,Koch2010}, and source and sink
edges are the inflow and outflow reactions of an open CRN. The signed
defect content $q(\sigma)=\chi n$ is a linear invariant (a left null
vector of $\mathsf{S}$) of every intra-host edge. Cluster dynamics also
descends from the Smoluchowski population balance
\cite{Smoluchowski1917}, itself an infinite CRN. The RAG is therefore a
\emph{descendant} of CRN theory in formalism, a \emph{sibling} of
rule-based network generators in method, and an \emph{extension} of both
in physics (Table~\ref{tab:S1_crn}).

The sibling relation guides the implementation. BioNetGen and Kappa
\cite{Faeder2009,DanosLaneve2004} generate reactions from rules over
structured species, and RMG \cite{Gao2016} organizes combustion mechanisms
into reaction families with rate rules. Likewise, an \texttt{Edge} is a
family over the size axis, its \texttt{EdgeClass} is the template, and the
kernel library $k$ supplies the rate rules. The RAG adds a typed state
space in which admissibility is checked structurally: (i)~polarity splits
$\mathcal{V}$ into two layers, annihilation being the only binary
cross-layer class and inter-population transfer being confined to one
layer (\texttt{Edge.\_\_post\_init\_\_}); (ii)~each of the ten edge
classes has a fixed stoichiometric contract (\texttt{EDGE\_CLASS\_SPEC}),
so conservation of $q$ is a property of the graph rather than of each
host's equations; (iii)~populations encode mobility and environment, so
one size can occupy distinct vertices (e.g., glissile
$\tfrac12\langle111\rangle$ and sessile $\langle100\rangle$ loops).

These extensions also mark where CRN theory stops applying. The
deficiency and complex-balancing theorems
\cite{HornJackson1972,Feinberg1987} assume finite networks with constant
rate coefficients, whereas RAG kernels are size-dependent,
diffusion-limited and tied to emission through binding energies;
$\mathcal{D}_a c_a$ is a mean-field closure; and the system is open and
driven far from equilibrium. The bin-moment and helium reductions coarsen
the graph while keeping the exact edge set, a graph-level analog of
lumping \cite{WeiKuo1969}. Finally, the cluster-dynamics ``master
equation'' is the deterministic rate equation~\eqref{eq:S1_walker}, not
the chemical master equation; the latter corresponds to stochastic
cluster dynamics \cite{Gillespie1976,MarianBulatov2011}, which could walk
the same edge set.

\begin{table}[t]
\centering
\caption{CRN concepts and their realization in the RAG and RadCluster.}
\label{tab:S1_crn}
\small
\begin{tabularx}{\textwidth}{@{}p{0.23\textwidth}p{0.30\textwidth}X@{}}
\toprule
CRN concept & RAG realization & Code (RadCluster\_2\_1/\texttt{py\_utils/core}) \\
\midrule
Species & Vertex $\sigma=(\chi,n,p,\mathbf{c})$ & \texttt{ClusterIdentifier}, \texttt{Population}, \texttt{Polarity} \\
Reaction & Edge $r$ with tuple $\nu_r$ & \texttt{Edge} (one \emph{family} over $n$) \\
Reaction template / rule & Edge class (10 classes) & \texttt{EdgeClass}, \texttt{EDGE\_CLASS\_SPEC} \\
Rate constant / rate rule & Kernel $k_r(n,n')$ & \texttt{register\_kernel} \\
Stoichiometric matrix $\mathsf{S}$ & $S_{ar}$, implied by the class contract & \texttt{GraphWalker} \\
Conservation law & $q=\chi n$ (and He inventory) & structural; checked by $\delta_{\rm FP}$, $\delta_{\rm He}$ \\
Inflow / outflow & \textsc{source} / \textsc{sink} classes & cascade spectrum; $\mathcal{D}_a$ \\
Species--reaction graph & Explicit hypergraph view & \texttt{to\_networkx} (diagnostics) \\
Model reduction (lumping) & Bin-moment and He reductions & \texttt{reductions/} \\
\emph{No CRN analog} & Polarity layers; population transfer; admissibility checks & \texttt{Edge.\_\_post\_init\_\_}, \texttt{rag.validate} \\
\bottomrule
\end{tabularx}
\end{table}

% ---------------------------------------------------------------------
\paragraph{Graph representation and use of network software.}
In production the RAG is implicit. \texttt{ReactionAdmissibilityGraph}
holds a population registry, a list of \texttt{Edge} families and a kernel
library, and rejects an edge whose endpoints are unregistered or whose
polarity rules fail. No vertex object is created: the walker applies each
family's kernel vectorized over size, so one edge stands for $O(N)$ unary
or $O(N^2)$ binary reactions. The 24 EUROFER-97 families
(Table~\ref{tab:S1_layer2}) would expand to $10^7$--$10^9$ arcs at
$I=V=10^3$--$10^4$, so the solver uses arrays, not a graph library.

For diagnostics, \texttt{to\_networkx} materializes the graph with
NetworkX \cite{Hagberg2008} as a \texttt{MultiDiGraph} truncated at
$n_{\max}$, with one node per admissible vertex and sentinel
\texttt{SOURCE} and \texttt{SINK} nodes. Unary families become
$n\to n\pm1$ or $(p,n)\to(p',n)$ arcs, and solute trapping a self-loop.
Because NetworkX has no hyperedge type, each binary edge is stored as its
two tail$\to$head arcs sharing a \texttt{hyper\_id}, which preserves the
hypergraph exactly. This representation serves three purposes.
(a)~\texttt{structural\_report} checks weak connectivity, reachability
from the monomer roots (an unreachable vertex is a population with no
production channel) and vertices without outgoing arcs. For EUROFER-97
at $n_{\max}=30$ it reports 89 vertices and 7078 arcs, full connectivity
and no unreachable or dead vertices; coalescence and annihilation account
for 93\% of the arcs, which is why production keeps the graph implicit.
(b)~\texttt{write\_graphml} exports to Gephi, yEd or Cytoscape.
(c)~The main-text RAG figures are drawn from this graph, not from a
hand-written edge list, so they cannot drift from the model. NetworkX is
imported lazily and is not required by the solver.

% ---------------------------------------------------------------------
\paragraph{Three-layer architecture.}
RadCluster separates what a cluster-dynamics model \emph{is}, what a
material \emph{admits}, and how the equations are \emph{integrated}
(Fig.~\ref{fig:S1_arch}). Following the open--closed principle
\cite{Meyer1988}, Layers~1 and~3 are closed to modification and a new
material enters only by extension in Layer~2. Both closed layers are
host-independent and verified once, so every material inherits their
conservation structure, reductions, error control and preconditioners;
editing either for one material would silently change all others.

\emph{Layer~1, the abstract core} (\texttt{py\_utils/core}), defines the
state-space vocabulary (\texttt{Polarity}, \texttt{Population},
\texttt{ClusterIdentifier}), the ten edge classes and their contracts, the
RAG container, the \texttt{StateLayout} map from vertices to the ODE
vector, the reference \texttt{GraphWalker} that assembles
Eq.~\eqref{eq:S1_walker} by dispatching on edge class, and the bin-moment
and helium reductions. It names no material.

\emph{Layer~2, the host declaration} (\texttt{materials/<host>} and its
compiled mirror in \texttt{cpp\_utils}), is the only code written for a
new material. It returns the pair (RAG, \texttt{StateLayout}) and installs
a compiled right-hand side for Layer~3. Its selections, with the EUROFER-97
choices in Table~\ref{tab:S1_layer2}, are:
\begin{enumerate}
\item \textbf{Populations}, each with $n_{\min}$, a mobility cut-off and a
  monomer pool per polarity layer. A separate population is warranted when
  clusters of equal size differ in mobility, binding energy or sink
  coupling (bcc $\tfrac12\langle111\rangle$ vs.\ $\langle100\rangle$ loops;
  hcp $\langle a\rangle$ vs.\ $\langle c\rangle$; fcc Frank loops vs.\
  stacking-fault tetrahedra).
\item \textbf{Composition axis}: the resolved non-host species (He for
  steels; He and H for plasma-facing W). An empty list removes solute
  trapping.
\item \textbf{Edge families}: each process mapped to a class, with partner
  and product populations. The core rejects inadmissible choices; excluded
  channels are simply not declared, and optional ones are switched by input
  flags (e.g., \texttt{LOOP\_NETWORK\_LOSS}).
\item \textbf{Kernel library $k$}: size-resolved rates from diffusivities,
  capture radii, bias factors, binding energies, sink strengths and the
  cascade spectrum, computed by the physics modules from the input
  workbook. $(\mathcal{V},\mathcal{E})$ is host-specific and $k$ is
  calibration-specific, so recalibration changes data, not topology.
\item \textbf{State representation}: discrete or bin-moment size axes and
  the helium reduction (Case~1 for fusion, Case~2 for fission), which fix
  the \texttt{StateLayout} and $N_{\rm eq}$.
\item \textbf{Compiled right-hand side} (\texttt{rate\_kernels.cpp}),
  verified against the Python \texttt{GraphWalker} in the order: walker,
  conservation gate ($\delta_{\rm FP},\delta_{\rm He}<10^{-6}$), C++
  mirror, calibration.
\end{enumerate}

\emph{Layer~3, the numerical engine} (\texttt{cpp\_utils/core},
\texttt{cpp\_bridge.py}), contains the CVODE variable-order BDF driver
\cite{Hindmarsh2005,Gardner2022}, which enforces non-negativity by a
post-step floor rather than solver constraints (these would create a
Jacobian kink); dense, banded, SPGMR \cite{SaadSchultz1986} (Jacobi or
Woodbury preconditioned) and KLU \cite{DavisKLU2010} linear solvers, the
last with a colored finite-difference Jacobian \cite{CPR1974}; the OpenMP
active window; and the Python bridge that launches the solver and parses
its output. It reaches the physics only through the function pointer
\texttt{UserData.rhs\_fn}.

\begin{figure}[p]
\centering
% Fig. S1 (architecture of RadCluster along its three layers).
% Needs: \usetikzlibrary{arrows.meta,positioning,calc}
\definecolor{Lone}{RGB}{52,86,139}     % Layer 1 accent
\definecolor{Ltwo}{RGB}{196,98,16}     % Layer 2 accent
\definecolor{Lthree}{RGB}{40,120,90}   % Layer 3 accent
\definecolor{Ldiag}{RGB}{105,105,105}
\resizebox{\textwidth}{!}{%
\begin{tikzpicture}[
  font=\sffamily\scriptsize,
  >={Stealth[length=2.2mm]},
  % b={color}{width}: a component box of given color and width (cm)
  b/.style 2 args={draw=#1, fill=#1!8, rounded corners=2pt, line width=0.6pt,
      align=center, inner sep=2pt, text width=#2 cm - 4pt,
      minimum width=#2 cm, minimum height=15mm},
  frame/.style={draw=#1, fill=#1!3, line width=1.1pt, rounded corners=5pt},
  lbl/.style={font=\sffamily\bfseries\normalsize, anchor=west, text=#1},
  det/.style={font=\sffamily\scriptsize, anchor=east, text=#1},
  note/.style={font=\sffamily\scriptsize\itshape, text=black!70},
  flow/.style={->, line width=1pt},
  pics/lock/.style={code={
      \draw[line width=1pt, #1] (-0.10,0.06) -- (-0.10,0.16)
            arc[start angle=180, end angle=0, radius=0.10] -- (0.10,0.06);
      \fill[#1] (-0.16,-0.17) rectangle (0.16,0.07);
      \fill[white] (0,-0.05) circle (0.035);
  }},
]
\def\W{15.9}   % total width

% ======================= LAYER 1: abstract core ==========================
\draw[frame=Lone] (-0.1,0.0) rectangle (\W+0.1,2.55);
\pic at (0.25,2.2) {lock=Lone};
\node[lbl=Lone] at (0.5,2.2) {Layer 1 --- Abstract core};
\node[det=Lone] at (\W,2.2)
  {\texttt{py\_utils/core}\quad host-independent $\cdot$ immutable};

\foreach \i/\name/\txt in {
  0/L1a/{\textbf{Cluster identifier}\\ $\sigma=(\chi,n,p,\mathbf c)$\\ {\tiny\texttt{Polarity}, \texttt{Population}}},
  1/L1b/{\textbf{10 edge classes}\\ stoichiometric contracts\\ {\tiny\texttt{EDGE\_CLASS\_SPEC}}},
  2/L1c/{\textbf{RAG container}\\ $\mathcal G=(\mathcal V,\mathcal E,\nu,k)$\\ admissibility checks},
  3/L1d/{\textbf{GraphWalker}\\ reference RHS,\\ Eq.~(\ref{eq:S1_walker})},
  4/L1e/{\textbf{StateLayout}\\ vertex $\leftrightarrow$ ODE index\\ (blocks)},
  5/L1f/{\textbf{Reductions}\\ bin moments, $P{=}2$\\ He Case 1/2}}
{
  \node[b={Lone}{2.15}, anchor=west] (\name) at ({\i*2.28},1.0) {\txt};
}
\node[b={Ldiag}{2.1}, fill=white, dashed, anchor=east] (L1g) at (\W,1.0)
  {\textbf{Diagnostics}\\ {\tiny\texttt{to\_networkx}}\\ NetworkX, optional,\\ off solver path};

% ======================= LAYER 2: host declaration =======================
\draw[frame=Ltwo, line width=1.5pt] (-0.1,-1.6) rectangle (\W+0.1,-7.0);
\node[lbl=Ltwo] at (0.1,-1.92) {Layer 2 --- Host declaration};
\node[det=Ltwo] at (\W,-1.92)
  {written per material};

% selection card
\node[draw=Ltwo, fill=Ltwo!10, rounded corners=2pt, line width=0.8pt,
      anchor=north west, align=left, text width=4.75cm, inner sep=4pt]
  (L2sel) at (0,-2.3)
  {\textbf{\normalsize\texttt{materials/eurofer97}}\\[3pt]
   \textbf{1}\; populations: {\tiny\texttt{bulk-111}, \texttt{bulk-100}, \texttt{bulk}}\\[1pt]
   \textbf{2}\; composition axis: He\\[1pt]
   \textbf{3}\; 24 edge families (P1--P8, loops)\\[1pt]
   \textbf{4}\; kernel library $k$ (calibrated)\\[1pt]
   \textbf{5}\; discrete/bin-moment; He Case 1/2\\[1pt]
   \textbf{6}\; compiled RHS mirror, verification};

\node[b={Ltwo}{2.6}, anchor=north west] (L2decl) at (5.25,-2.35)
  {\textbf{declaration}\\ {\tiny\texttt{build\_eurofer\_rag()}}\\ $\to(\mathcal G_{\rm Eur}$, {\tiny\texttt{StateLayout}}$)$};
\node[b={Ltwo}{2.45}, anchor=north west] (L2phys) at (8.55,-2.35)
  {\textbf{Physics modules}\\ {\tiny\texttt{reaction\_rates}\\ \texttt{binding\_energies}\\ \texttt{defect\_production}}};
\node[b={gray}{1.75}, fill=black!4, anchor=north west] (inp) at (11.3,-2.35)
  {\textbf{Input}\\ workbook,\\ calibration};
\node[b={Ltwo}{2.6}, anchor=north west] (L2cpp) at (5.25,-4.75)
  {\textbf{compiled RHS}\\ {\tiny\texttt{rate\_kernels.cpp}}\\ $\to$ {\tiny\texttt{rhs\_fn}}};

% other hosts: extension points
\foreach \j/\name/\txt in {
  0/H1/{{\tiny\texttt{materials/zr}}\\ hcp: $\langle a\rangle$, $\langle c\rangle$ families},
  1/H2/{{\tiny\texttt{materials/w}}\\ composition He, H},
  2/H3/{{\tiny\texttt{materials/<fcc>}}\\ Frank loops, SFTs}}
{
  \node[b={Ltwo}{2.6}, fill=white, dashed, minimum height=12mm, anchor=north east]
    (\name) at (\W,{-2.35-\j*1.45}) {\txt};
}
\node[note, anchor=north east] at (\W,-6.55) {further hosts reuse Layers 1 and 3};

\draw[flow, Ltwo] (inp) -- (L2phys);
\draw[flow, Ltwo] (L2phys) -- (L2decl);
\draw[flow, Ltwo] (L2sel.east |- L2decl) -- (L2decl.west);
\draw[flow, Ltwo, densely dashed] (L2decl) -- node[left, note]{mirror} (L2cpp);

% ======================= LAYER 3: numerical engine =======================
\draw[frame=Lthree] (-0.1,-8.2) rectangle (\W+0.1,-10.75);
\pic at (0.25,-8.55) {lock=Lthree};
\node[lbl=Lthree] at (0.5,-8.55) {Layer 3 --- Numerical engine};
\node[det=Lthree] at (\W,-8.55)
  {\texttt{cpp\_utils/core}, \texttt{cpp\_bridge.py}\quad host-independent $\cdot$ immutable};

\foreach \i/\name/\txt in {
  0/L3a/{\textbf{Python bridge}\\ parameter file,\\ launch, parse output},
  1/L3b/{\textbf{CVODE (BDF)}\\ tolerances, $C_{\rm floor}$,\\ segments},
  2/L3c/{\textbf{Linear solvers}\\ dense $\cdot$ band\\ SPGMR $\cdot$ KLU},
  3/L3d/{\textbf{Preconditioners}\\ Jacobi; Woodbury\\ $\mathsf T+\mathsf U\mathsf V^{\mathsf T}$},
  4/L3e/{\textbf{Sparse Jacobian}\\ CPR coloring,\\ FD fill},
  5/L3f/{\textbf{Active window}\\ sliding fronts,\\ OpenMP RHS}}
{
  \node[b={Lthree}{2.5}, anchor=west] (\name) at ({\i*2.68},-9.7) {\txt};
}

% outputs
\node[b={gray}{9.0}, fill=black!4, minimum height=8mm] (out) at (\W/2,-11.65)
  {trajectories $c_a(t)$ $\cdot$ swelling, loop and void statistics $\cdot$
   $\delta_{\rm FP}$, $\delta_{\rm He}$ $\cdot$ \texttt{provenance.md}};
\draw[flow, Lthree] (\W/2,-10.75) -- (out.north);

% ======================= inter-layer flows ===============================
\draw[flow, Lone, line width=1.3pt] (L1c.south) -- (L1c.south |- L2decl.north)
  node[pos=0.55, left, note, align=right]{instantiate:\\ types, contracts};
\draw[flow, Ltwo, line width=1.5pt] (L2cpp.south) -- (L2cpp.south |- 0,-8.2)
  node[pos=0.62, right, note, align=left]
  {\texttt{rhs\_fn}, \texttt{StateLayout},\\ parameters $(i_{\rm mobile}, v_{\rm mobile})$};
\draw[<->, line width=0.9pt, densely dashed, Lone]
  (L1d.south) -- ++(0,-0.95) -| (8.2,-5.5) -- (L2cpp.east |- 0,-5.5)
;
\node[note, anchor=north west, align=left] at (8.35,-4.3)
  {verify: walker $=$ compiled RHS\\ (conservation gate)};
\end{tikzpicture}}

\caption{Three-layer architecture of RadCluster. Layers~1 and~3 are
host-independent and closed to modification (padlocks); a material enters
only through Layer~2, which passes Layer~3 a \texttt{StateLayout} and a
compiled right-hand side verified against the Layer~1 walker. NetworkX
diagnostics lie outside the solver path; dashed cards are further hosts.}
\label{fig:S1_arch}
\end{figure}

\begin{table}[t]
\centering
\caption{Layer-2 selections for EUROFER-97 in RadCluster\_2\_1
(\texttt{materials/eurofer97/declaration.py}).}
\label{tab:S1_layer2}
\small
\begin{tabularx}{\textwidth}{@{}p{0.20\textwidth}X@{}}
\toprule
Selection & EUROFER-97 instantiation \\
\midrule
SIA populations & \texttt{bulk-111} ($n_{\min}=1$, mobile to $i_{\rm mobile}$, SIA monomer pool, all cascade SIA production);
                  \texttt{bulk-100} (sessile, $n_{\min}=n_{\rm loop}$, default 4) \\
Vacancy populations & \texttt{bulk} ($n_{\min}=1$, mobile to $v_{\rm mobile}$, vacancy monomer pool) \\
Composition axis & He only \\
Edge families (24) & P1 recombination; P2 cavity growth/shrinkage; P3 loop growth/shrinkage (both SIA populations);
  P4 fixed sinks (3); P5 thermal emission (3); SIA--SIA and V--V coalescence; V--I annihilation with
  $\tfrac12\langle111\rangle$ and $\langle100\rangle$; $\tfrac12\langle111\rangle\to\langle100\rangle$
  conversion by unary transfer (Dudarev), junction coalescence and absorption growth (Marian);
  cascade sources (2); P7 trap mutation; P8 radiation re-solution \\
Excluded channels & 1D/1D reactions between mutually mobile clusters \\
Kernel inputs & $\Omega=a^3/2$; $D_n$ with solute trapping; $\xi(n)=(3n/8\pi)^{1/3}$; bias factors; $E_b(n)$; $\rho_d$, $d_g$, precipitates; cascade spectra \\
State representation & discrete ($N_{\rm eq}\approx 2I+V$) or bin-moment (linear, $P=2$); He Case~2 (fission) or Case~1 (fusion); five conservation integrals \\
\bottomrule
\end{tabularx}
\end{table}

% ---------------------------------------------------------------------
\paragraph{Relationship between the solver method and the RAG.}
The RAG is a model and the integrator a method: the graph fixes the vector
field of Eq.~\eqref{eq:S1_walker} and the structure of its Jacobian
$\mathsf{J}_{ab}=\partial f_a/\partial c_b$, but not how the equations are
advanced. The same edge set can be walked by the Python reference, by
compiled kernels under CVODE, by the bin-moment
reconstruct--walk--project wrapper or by a stochastic algorithm. The
graph does, however, determine which method is efficient.

\emph{(i) Jacobian sparsity is RAG adjacency.} Because $J_r$ depends only
on the tail vertices of $r$, $\mathsf{J}_{ab}\neq0$ only if some edge has
$b$ in its tail and $a$ in its tail or head (Fig.~\ref{fig:S1_jac}).
Unary ladder edges couple $n$ to $n\pm1$. Each monomer lies in the tail of
$O(N)$ reactions, giving a dense column. Binary edges with a mobile
partner $m\le i_{\rm mobile}$ couple $n$ to $n+m$, widening the band and
making all mobile columns dense, while sessile--sessile pairs have zero
kernel. Hence
\begin{equation}
  \mathsf{J} \;\approx\; \mathsf{T} + \mathsf{U}\mathsf{V}^{\mathsf T},\qquad
  b_w=\max(2i_{\rm mobile},2v_{\rm mobile})+1,\qquad
  r=\operatorname{rank}\mathsf{U}=i_{\rm mobile}+v_{\rm mobile},
  \label{eq:S1_woodbury}
\end{equation}
with $\mathsf T$ banded (half-width $b_w$) and $\mathsf{V}$ selecting the
mobile columns. These are the defaults of the Woodbury preconditioner
(\texttt{prec\_bw}, \texttt{prec\_rank}): the band is probed with $2b_w+1$
colored evaluations and factored with LAPACK, the $r$ mobile columns are
probed individually, and $(\mathsf I-\gamma\mathsf J)\mathbf x=\mathbf b$
is solved by the Sherman--Morrison--Woodbury identity \cite{Hager1989}
with an $r\times r$ Schur complement. Long-range entries in the mobile
rows are dropped, an approximation GMRES tolerates. The KLU path uses the
same pattern explicitly, and its color count sets the evaluations per
Jacobian.

\emph{(ii) Stiffness is the spread of edge rates.} Monomer recombination
and small-cluster migration are many orders of magnitude faster than
absorption by large sessile clusters, hence implicit BDF integration. For
discrete systems ($N_{\rm eq}\sim10^4$), preconditioned GMRES outperforms
sparse direct factorization, whose colored Jacobian rebuilds dominate the
cost; for bin-moment systems ($N_{\rm eq}\sim10^2$), per-size
reconstruction dominates and the preconditioner matters less.

\emph{(iii) Reachability governs the active window.} Vertices beyond the
growth front are reachable only through growth or coalescence edges
leaving it. The active window integrates this reachable subgraph and
extends it by a fixed pad when the front concentration exceeds a
threshold, a solution-driven breadth-first expansion from the monomer
roots.

\emph{(iv) Separating topology from kernels preserves solver structure.}
Slowly evolving kernels can be handled by operator splitting: the network
dislocation density of the loop-to-network loss channel is frozen within a
segment and the kernels rebuilt between segments (\texttt{rebuild\_rates})
with the edge set unchanged, so the sparsity pattern, coloring and
preconditioner carry over. Bin moments replace vertex ranges by
super-vertices yet evaluate exact per-size fluxes, inheriting
conservation; because they alter the coupling pattern, the colored-KLU
path is restricted to discrete modes. Since conservation of $q$ is built
into the class contracts, $\delta_{\rm FP}$ and $\delta_{\rm He}$ measure
integration error alone.

\begin{figure}[t]
\centering
% Fig. S1 (RAG -> Jacobian structure).  Needs tikz libraries:
%   arrows.meta, positioning, calc, decorations.pathreplacing, patterns
\definecolor{Jband}{RGB}{52,86,139}
\definecolor{Jdense}{RGB}{196,98,16}
\definecolor{Jaux}{RGB}{40,120,90}
\resizebox{0.92\textwidth}{!}{%
\begin{tikzpicture}[font=\sffamily\footnotesize, >={Stealth[length=2mm]}]

% ---------------- left: a size ladder of the RAG ------------------------
\begin{scope}[shift={(0,0)}]
  \node[font=\sffamily\bfseries\small, anchor=west] at (-0.3,3.55) {(a) RAG, one size ladder};
  % mobile vertices (hubs)
  \foreach \n/\x in {1/0,2/0.9,3/1.8}{
    \node[circle, draw=Jdense, fill=Jdense!25, line width=0.8pt,
          minimum size=6.5mm, inner sep=0pt] (m\n) at (\x,2.3) {\n};
  }
  % sessile vertices
  \foreach \n/\x in {4/2.7,5/3.6,6/4.5,7/5.4}{
    \node[circle, draw=Jband, fill=Jband!12, line width=0.8pt,
          minimum size=6.5mm, inner sep=0pt] (s\n) at (\x,0.6) {\n};
  }
  \node at (6.25,0.6) {$\cdots$};
  \node[circle, draw=Jband, fill=Jband!12, line width=0.8pt,
        minimum size=6.5mm, inner sep=0pt] (sN) at (7.0,0.6) {$N$};
  % unary ladder edges n -> n+1 among sessile
  \foreach \a/\b in {s4/s5,s5/s6,s6/s7}{
    \draw[->, Jband, line width=0.7pt] (\a) to[bend left=35] (\b);
    \draw[->, Jband!60, line width=0.5pt] (\b) to[bend left=35] (\a);
  }
  % hub arcs: monomer (1) participates in every ladder edge
  \foreach \t in {s4,s5,s6,s7,sN}{
    \draw[Jdense, line width=0.45pt, densely dashed] (m1) -- (\t);
  }
  % binary hyperedge: m2 + s4 -> s6
  \coordinate (hx) at (3.55,1.55);
  \draw[Jdense!80!black, line width=0.8pt] (m2) -- (hx);
  \draw[Jdense!80!black, line width=0.8pt] (s4) -- (hx);
  \draw[->, Jdense!80!black, line width=0.8pt] (hx) -- (s6);
  \fill[Jdense!80!black] (hx) circle (1.3pt);
  \node[font=\sffamily\scriptsize, Jdense!80!black, anchor=west] at (3.7,1.75)
       {$I_2+I_4\to I_6$};
  \draw[decorate, decoration={brace, amplitude=4pt}, line width=0.6pt]
       (-0.35,2.75) -- (2.15,2.75)
       node[midway, above=3pt, font=\sffamily\scriptsize]{mobile ($n\le i_{\rm mobile}$): hubs};
  \node[font=\sffamily\scriptsize, Jband, anchor=north] at (4.6,0.05)
       {sessile: ladder edges; sessile--sessile kernels $=0$};
\end{scope}

% ---------------- right: Jacobian pattern -------------------------------
\begin{scope}[shift={(9.2,-0.9)}]
  \def\L{4.6}     % matrix side
  \def\Is{2.6}    % SIA block size
  \def\r{0.28}    % mobile strip width (SIA)
  \def\rv{0.16}   % mobile strip width (VAC)
  \def\bw{0.30}   % half band
  \def\ax{0.22}   % aux rows
  \node[font=\sffamily\bfseries\small, anchor=west] at (-0.3,\L+0.85) {(b) Jacobian $\mathsf J_{ab}=\partial f_a/\partial c_b$};
  % band T along diagonal of each block (drawn top-left to bottom-right)
  \fill[Jband!30] (0,\L) -- (\bw,\L) -- (\Is,\L-\Is+\bw) -- (\Is,\L-\Is)
        -- (\Is-\bw,\L-\Is) -- (0,\L-\bw) -- cycle;
  \fill[Jband!30] (\Is,\L-\Is) -- (\Is+\bw,\L-\Is) -- (\L,\ax+\bw) -- (\L,\ax)
        -- (\L-\bw,\ax) -- (\Is,\L-\Is-\bw) -- cycle;
  % dense mobile columns (U V^T)
  \fill[Jdense!55] (0,\ax) rectangle (\r,\L);
  \fill[Jdense!55] (\Is,\ax) rectangle (\Is+\rv,\L);
  % dense mobile rows (omitted by Woodbury outside the band)
  \fill[pattern=north east lines, pattern color=Jdense] (0,\L-\r) rectangle (\L,\L);
  \fill[pattern=north east lines, pattern color=Jdense] (0,\L-\Is-\rv) rectangle (\L,\L-\Is);
  % aux (He, conservation) rows
  \fill[Jaux!45] (0,0) rectangle (\L,\ax);
  % frame and block separators
  \draw[line width=0.8pt] (0,0) rectangle (\L,\L);
  \draw[black!50, densely dotted] (\Is,0) -- (\Is,\L);
  \draw[black!50, densely dotted] (0,\L-\Is) -- (\L,\L-\Is);
  \node[font=\sffamily\scriptsize, anchor=south] at (\Is/2,\L) {SIA};
  \node[font=\sffamily\scriptsize, anchor=south] at ({(\Is+\L)/2},\L) {VAC};
  \node[font=\sffamily\scriptsize, anchor=east] at (0,\ax/2) {He, cons.};
  % band width annotation
  \draw[<->, line width=0.5pt] (1.3,\L-1.3) -- ++(\bw*0.71,\bw*0.71)
        node[right, font=\sffamily\scriptsize, fill=white, inner sep=1pt] {$b_w$};
  % legend
  \begin{scope}[shift={(-9.0,-0.35)}, font=\sffamily\footnotesize]
    \fill[Jband!30] (0,0) rectangle (0.35,0.25);
    \node[anchor=west] at (0.45,0.12) {band $\mathsf T$: ladder + mobile-partner reach};
    \fill[Jdense!55] (0,-0.5) rectangle (0.35,-0.25);
    \node[anchor=west] at (0.45,-0.38) {dense columns $\mathsf U\mathsf V^{\mathsf T}$, rank $r$};
    \fill[pattern=north east lines, pattern color=Jdense] (0,-1.0) rectangle (0.35,-0.75);
    \node[anchor=west] at (0.45,-0.88) {dense mobile rows (not in preconditioner)};
    \fill[Jaux!45] (0,-1.5) rectangle (0.35,-1.25);
    \node[anchor=west] at (0.45,-1.38) {aggregate He / conservation rows};
  \end{scope}
  \node[anchor=north, align=center, font=\sffamily\footnotesize] at (\L/2,-0.15)
      {$b_w=\max(2i_{\rm mobile},2v_{\rm mobile})+1$,\quad
       $r=i_{\rm mobile}+v_{\rm mobile}$};
\end{scope}

\draw[->, line width=1.1pt, black!60] (7.6,1.4) -- node[above, font=\sffamily\scriptsize]{$b\in\mathrm{tail}(r)$} (8.8,1.4);
\end{tikzpicture}}

\caption{RAG connectivity and the Jacobian pattern of the SIA and vacancy
ladders. Ladder edges give the tridiagonal core, binary edges with mobile
partners widen it to the band $\mathsf T$, and mobile hub vertices give
the dense columns $\mathsf U\mathsf V^{\mathsf T}$ and dense rows. The
Woodbury preconditioner retains $\mathsf T$ and the mobile columns.}
\label{fig:S1_jac}
\end{figure}
\endgroup
